\documentclass{beamer}

\usetheme{Madrid}
\usecolortheme{seahorse}
\usepackage[utf8]{inputenc}
\usepackage[T1]{fontenc}
\usepackage{listings}
\usepackage{xcolor}

% --- Configuración del estilo del código C# ---
\lstdefinestyle{CSharpStyle}{
  language=[Sharp]C,
  basicstyle=\scriptsize\ttfamily,
  keywordstyle=\color{blue}\bfseries,
  stringstyle=\color{teal},
  commentstyle=\color{gray}\itshape,
  showstringspaces=false,
  breaklines=true,
  frame=single,
  rulecolor=\color{gray!30},
  tabsize=2,
  captionpos=b,
  numbers=left,
  numberstyle=\tiny\color{gray},
}

% --- Título ---
\title{Uso de la Cola (Queue) en C\# para Videojuegos en Unity}
\author{Luis Antonio Molina Yampa}
\institute{COM200 - Programacion Avanzada}
\date{\today}

\begin{document}

% --- Portada ---
\begin{frame}
  \titlepage
\end{frame}

% --- Concepto General ---
\begin{frame}{Concepto General}
  Una \textbf{Queue (cola)} en C\# es una estructura de datos \textbf{FIFO (First In, First Out)}, es decir, el primer elemento en ingresar es el primero en salir.  

  \vspace{0.3cm}
  En el contexto de Unity, se puede usar para gestionar \textbf{enemigos que aparecen en oleadas}, asegurando que se procesen en el orden correcto.  

  \vspace{0.3cm}
  Ejemplos de uso:
  \begin{itemize}
    \item Control de turnos de enemigos o jugadores.
    \item Spawns secuenciales en escenas.
    \item Ejecución ordenada de eventos o tareas.
  \end{itemize}
\end{frame}

% --- Declaración de una Queue ---
\begin{frame}[fragile]{Declaración de una Queue en C\#}
\lstset{style=CSharpStyle}
Antes de usar una cola, se importa:
\begin{lstlisting}
using System.Collections.Generic;
\end{lstlisting}

Se puede declarar una cola de objetos que representen enemigos:
\begin{lstlisting}
Queue<GameObject> colaDeEnemigos = new Queue<GameObject>();
\end{lstlisting}
\end{frame}

% --- Agregar elementos a la cola ---
\begin{frame}[fragile]{Agregar enemigos a la cola — Enqueue()}
\lstset{style=CSharpStyle}
El método \textbf{Enqueue()} agrega un elemento al final de la cola:
\begin{lstlisting}
colaDeEnemigos.Enqueue(goblin);
colaDeEnemigos.Enqueue(orco);
colaDeEnemigos.Enqueue(dragon);
\end{lstlisting}

\textbf{Ejemplo práctico:}  
Cuando un enemigo entra al mapa, se añade a la cola para controlar el orden de aparición.
\end{frame}

% --- Quitar el primer elemento ---
\begin{frame}[fragile]{Procesar enemigos — Dequeue()}
\lstset{style=CSharpStyle}
\textbf{Dequeue()} elimina y devuelve el primer enemigo de la cola:
\begin{lstlisting}
GameObject siguiente = colaDeEnemigos.Dequeue();
Debug.Log("Atacando al enemigo: " + siguiente.name);
\end{lstlisting}

Se recomienda verificar si la cola tiene elementos antes de llamar a \textbf{Dequeue()}:
\begin{lstlisting}
if(colaDeEnemigos.Count > 0)
{
    GameObject enemigoActual = colaDeEnemigos.Dequeue();
}
\end{lstlisting}
\end{frame}

% --- Ver el siguiente elemento ---
\begin{frame}[fragile]{Consultar sin eliminar — Peek()}
\lstset{style=CSharpStyle}
El método \textbf{Peek()} permite ver el primer enemigo sin retirarlo:
\begin{lstlisting}
GameObject siguiente = colaDeEnemigos.Peek();
Debug.Log("El siguiente enemigo será: " + siguiente.name);
\end{lstlisting}

En Unity, esto puede utilizarse para mostrar en la UI quién será el próximo en actuar.
\end{frame}

% --- Vaciar la cola y contar elementos ---
\begin{frame}[fragile]{Otras operaciones importantes}
\begin{itemize}
    \item \textbf{Clear()} — Vacía toda la cola.
\begin{lstlisting}
colaDeEnemigos.Clear();
\end{lstlisting}

    \item \textbf{Count} — Devuelve la cantidad de enemigos en la cola.
\begin{lstlisting}
int cantidad = colaDeEnemigos.Count;
\end{lstlisting}

    \item \textbf{Contains()} — Verifica si un enemigo está en la cola.
\begin{lstlisting}
if(colaDeEnemigos.Contains(dragon))
    Debug.Log("El dragón está en espera");
\end{lstlisting}
\end{itemize}
\end{frame}

% --- Ejemplo completo en Unity ---
\begin{frame}[fragile]{Ejemplo completo: Spawns de enemigos}
\lstset{style=CSharpStyle}
\begin{lstlisting}
using UnityEngine;
using System.Collections.Generic;

public class SpawnManager : MonoBehaviour
{
    public GameObject goblin, orco, dragon;
    private Queue<GameObject> colaDeEnemigos = new Queue<GameObject>();

    void Start()
    {
        colaDeEnemigos.Enqueue(goblin);
        colaDeEnemigos.Enqueue(orco);
        colaDeEnemigos.Enqueue(dragon);

        StartCoroutine(SpawnearEnemigos());
    }

    System.Collections.IEnumerator SpawnearEnemigos()
    {
        while(colaDeEnemigos.Count > 0)
        {
            GameObject siguiente = colaDeEnemigos.Dequeue();
            Instantiate(siguiente, transform.position, Quaternion.identity);
            Debug.Log("Enemigo spawneado: " + siguiente.name);
            yield return new WaitForSeconds(2f);
        }
    }
}
\end{lstlisting}
\end{frame}

% --- Ventajas de Queue en Unity ---
\begin{frame}{Ventajas de usar Queue en Unity}
\begin{itemize}
    \item Garantiza que los enemigos o eventos se procesen en \textbf{orden lógico}.
    \item Simplifica la administración de oleadas o turnos.
    \item Evita que los objetos se procesen de manera desordenada.
    \item Permite integrar sistemas de UI mostrando al próximo enemigo en cola.
\end{itemize}
\end{frame}

% --- Comparación con otras estructuras ---
\begin{frame}{Comparación con Stack y List}
\begin{table}[]
\centering
\begin{tabular}{|l|l|l|}
\hline
\textbf{Estructura} & \textbf{Orden de acceso} & \textbf{Ejemplo en Unity} \\ \hline
Queue & FIFO (primero entra, primero sale) & Spawns de enemigos, turnos \\ \hline
Stack & LIFO (último entra, primero sale) & Recolector de platos, historial \\ \hline
List & Acceso por índice & Inventario, lista de objetos activos \\ \hline
\end{tabular}
\end{table}
\end{frame}

% --- Conclusión ---
\begin{frame}{Conclusión}
\begin{itemize}
    \item La \textbf{Queue} en C\# permite controlar elementos en orden secuencial.
    \item Es ideal para sistemas de spawns, turnos y tareas que deben procesarse paso a paso.
    \item Facilita la programación de escenas donde el orden de aparición o ejecución es crítico.
\end{itemize}
\end{frame}

\end{document}
